% ------------------------------------------------------------
% DATOS DEL LIBRO
% ------------------------------------------------------------

\newcommand{\tituloobra}{Física}
\newcommand{\autorlibro}{Julian David Chavarro Balanta}

\providecommand{\titulovolumen}{}
\providecommand{\subtitulovolumen}{}


% ------------------------------------------------------------
% MATEMÁTICAS GENERALES
% ------------------------------------------------------------

\newcommand{\abs}[1]{\left|#1\right|}
\newcommand{\norm}[1]{\left\lVert#1\right\rVert}

\newcommand{\fr}[2]{\dfrac{#1}{#2}}

\newcommand{\paren}[1]{\left(#1\right)}
\newcommand{\bracket}[1]{\left[#1\right]}
\newcommand{\braces}[1]{\left\{#1\right\}}


% ------------------------------------------------------------
% VECTORES
% ------------------------------------------------------------

\newcommand{\vect}[1]{\bm{#1}}

\newcommand{\uvec}[1]{\hat{\bm{#1}}}

\newcommand{\ihat}{\hat{\bm{\imath}}}
\newcommand{\jhat}{\hat{\bm{\jmath}}}
\newcommand{\khat}{\hat{\bm{k}}}

\newcommand{\pdot}{\mathbin{\bm{\cdot}}}
\newcommand{\pcross}{\mathbin{\bm{\times}}}


% ------------------------------------------------------------
% CÁLCULO
% ------------------------------------------------------------

\newcommand{\dd}{\mathop{}\!\mathrm{d}}

\newcommand{\deriv}[2]{
    \frac{\dd #1}{\dd #2}
}

\newcommand{\derivdos}[2]{
    \frac{\dd^{2}#1}{\dd #2^{2}}
}

\newcommand{\derivn}[3]{
    \frac{\dd^{#3}#1}{\dd #2^{#3}}
}

\newcommand{\pderiv}[2]{
    \frac{\partial #1}{\partial #2}
}

\newcommand{\pderivdos}[2]{
    \frac{\partial^{2}#1}{\partial #2^{2}}
}


% ------------------------------------------------------------
% SÍMBOLOS GENERALES
% ------------------------------------------------------------

\newcommand{\grado}{^{\circ}}
\newcommand{\implica}{\Longrightarrow}
\newcommand{\sii}{\Longleftrightarrow}
\newcommand{\proporcional}{\propto}


% ------------------------------------------------------------
% OPERADORES
% ------------------------------------------------------------

\DeclareMathOperator{\sen}{sen}
\DeclareMathOperator{\arcsen}{arcsen}


% ------------------------------------------------------------
% TABLAS
% ------------------------------------------------------------

\newcolumntype{Y}{>{\centering\arraybackslash}X}


% ------------------------------------------------------------
% NIVELES DE EJERCICIOS
% ------------------------------------------------------------

\newcommand{\nivelbasico}{
    \subsubsection*{Nivel básico}
}

\newcommand{\nivelintermedio}{
    \subsubsection*{Nivel intermedio}
}

\newcommand{\nivelavanzado}{
    \subsubsection*{Nivel avanzado}
}